\chapter{Complexity — Bidirectional}
%%% generated by the blueprint pipeline from TCSlib.Complexity.TuringMachine.Robustness.Bidirectional
\section{Overview}

This module formalizes the robustness of the multi-tape Turing machine model under
the restriction to unidirectional (one-way infinite) work tapes, following
Arora--Barak, Claim~1.8.  Since the project's tapes are $\bbz$-indexed, the claim is
rendered as: every machine computing a function within a time bound is simulated,
with the same number of work tapes and constant-factor slowdown, by a machine over an
enlarged alphabet whose work heads never visit a negative cell.  The simulator folds
each work tape at the origin, storing the payloads of a nonnegative cell and of its
mirror negative cell together in one physical cell, alongside a Boolean origin flag
written during a one-step initialization phase.  The chapter records the folding
coordinate arithmetic, the translated read/write/move operations with their
correctness lemmas, the step-by-step simulation invariants, and the safety invariant
guaranteeing nonnegative head positions on arbitrary enlarged-alphabet inputs.

\section{Declarations}

\begin{definition}[Unidirectional tape use]
\lean{Turing.FinTM.NonnegativeHeads}
\label{Turing.FinTM.NonnegativeHeads}
\leanok
\uses{Turing.FinTM, Turing.MultiTapeTM.initCfg, Turing.MultiTapeTM.runFrom}
Let $M$ be a multi-tape Turing machine over a finite alphabet $\Gamma$.  The
predicate asserts that $M$ uses its work tapes unidirectionally: for every input
string over $\Gamma$, every number of steps $t \in \bbn$, and every work tape of $M$,
the position of that work tape's head after running $M$ for $t$ steps from its
initial configuration on the input is nonnegative.
\end{definition}

\begin{definition}[Enlarged folded alphabet]
\lean{Turing.FinTM.FoldSymbol}
\label{Turing.FinTM.FoldSymbol}
\leanok
The enlarged tape alphabet used to fold a two-way infinite tape over a base alphabet
$\Gamma$: its symbols are triples $(b, u, v)$ where $b$ is a Boolean origin flag and
$u$ and $v$ are two independent payloads, each either a symbol of $\Gamma$ or blank.
The intent is that a nonnegative physical cell $p$ carries the content $u$ of the
simulated cell $p$ and the content $v$ of the simulated cell $-p-1$.
\end{definition}

\begin{definition}[Embedding of the base alphabet]
\lean{Turing.FinTM.foldEmbedding}
\label{Turing.FinTM.foldEmbedding}
\leanok
\uses{Turing.FinTM.FoldSymbol}
The injection of a base alphabet $\Gamma$ into the folded alphabet sending a symbol
$a$ to the triple whose origin flag is unset, whose first payload is $a$, and whose
second payload is blank.
\end{definition}

\begin{definition}[Folded head coordinate]
\lean{Turing.FinTM.foldPos}
\label{Turing.FinTM.foldPos}
\leanok
The coordinate-folding map $\varphi \colon \bbz \to \bbz$ defined by
$\varphi(z) = z$ when $z \ge 0$ and $\varphi(z) = -z-1$ otherwise.  It maps each
negative cell $z$ onto the nonnegative cell $-z-1$; in particular
$\varphi(0) = \varphi(-1) = 0$, so $\varphi$ is not the absolute value.
\end{definition}

\begin{definition}[Side of an integer coordinate]
\lean{Turing.FinTM.foldSide}
\label{Turing.FinTM.foldSide}
\leanok
The Boolean recording on which side of the fold an integer coordinate $z$ lies:
true when $z \ge 0$ and false when $z < 0$.
\end{definition}

\begin{definition}[Canonical packing of a folded symbol]
\lean{Turing.FinTM.foldPack}
\label{Turing.FinTM.foldPack}
\leanok
\uses{Turing.FinTM.FoldSymbol}
The canonical packing of a folded-alphabet triple into an optional cell content: the
untagged triple with both payloads blank is sent to the physical blank, and every
other triple is kept as an actual symbol.  This ensures that a physical cell never
stores a symbol carrying no information.
\end{definition}

\begin{definition}[Unpacking of a physical cell]
\lean{Turing.FinTM.foldUnpack}
\label{Turing.FinTM.foldUnpack}
\leanok
\uses{Turing.FinTM.FoldSymbol}
The unpacking of a physical cell content over the folded alphabet --- either a folded
triple or blank --- into a folded triple, interpreting a physical blank as the
untagged triple with both payloads blank.
\end{definition}

\begin{lemma}[Unpacking inverts packing]
\lean{Turing.FinTM.foldUnpack_pack}
\label{Turing.FinTM.foldUnpack_pack}
\leanok
\uses{Turing.FinTM.FoldSymbol, Turing.FinTM.foldPack, Turing.FinTM.foldUnpack}
\difficulty{2}
For every folded-alphabet triple $v$ --- a Boolean flag together with two payloads,
each a base symbol or blank --- unpacking the canonical packing of $v$ returns $v$
itself.
\end{lemma}

\begin{definition}[Folding of a work tape]
\lean{Turing.FinTM.foldTape}
\label{Turing.FinTM.foldTape}
\leanok
\uses{Turing.FinTM.FoldSymbol, Turing.FinTM.foldPack}
Given a two-way infinite tape content $t \colon \bbz \to \Gamma \cup \{\square\}$
over a base alphabet $\Gamma$ (with $\square$ the blank), the folded tape assigns to
each nonnegative physical cell $p$ the canonical packing of the triple consisting of
the flag ``$p = 0$'', the content $t(p)$, and the content $t(-p-1)$; every negative
physical cell is blank.
\end{definition}

\begin{definition}[Translation of a head move across the fold]
\lean{Turing.FinTM.foldMove}
\label{Turing.FinTM.foldMove}
\leanok
Given the current side of the fold, an origin flag, and a simulated move direction
$d \in \{-1, 0, +1\}$, this operation returns the physical move to perform together
with the new side: moves translate directly on the nonnegative side and are reversed
on the negative side, while a crossing between simulated cells $0$ and $-1$
(detected by the origin flag) flips the side without issuing a physical move.
\end{definition}

\begin{lemma}[Folded coordinates are nonnegative]
\lean{Turing.FinTM.foldPos_nonneg}
\label{Turing.FinTM.foldPos_nonneg}
\leanok
\uses{Turing.FinTM.foldPos}
\difficulty{2}
For every integer $z$, the folded coordinate --- equal to $z$ when $z \ge 0$ and to
$-z-1$ otherwise --- is nonnegative.
\end{lemma}

\begin{lemma}[Correctness of the translated move]
\lean{Turing.FinTM.foldMove_correct}
\label{Turing.FinTM.foldMove_correct}
\leanok
\uses{Turing.FinTM.foldMove, Turing.FinTM.foldPos, Turing.FinTM.foldSide}
\difficulty{4}
Let $z$ be an integer and $d \in \{-1,0,+1\}$ a move direction; write $\varphi(z)$
for the folded coordinate ($z$ when $z \ge 0$, else $-z-1$) and $\sigma(z)$ for the
side of $z$ (whether $z \ge 0$).  Then the translated move computed from $\sigma(z)$,
the flag ``$\varphi(z) = 0$'', and $d$ is correct in both components: adding its
physical displacement to $\varphi(z)$ gives $\varphi(z+d)$, and its returned side
equals $\sigma(z+d)$.
\end{lemma}

\begin{definition}[Reading the active payload]
\lean{Turing.FinTM.foldRead}
\label{Turing.FinTM.foldRead}
\leanok
\uses{Turing.FinTM.FoldSymbol, Turing.FinTM.foldUnpack}
Reading a physical cell over the folded alphabet relative to a side: after unpacking
the cell (a physical blank unpacks to the untagged triple of blanks), return the
first payload when the side is the nonnegative one and the second payload otherwise.
\end{definition}

\begin{lemma}[Folded reads agree with the virtual tape]
\lean{Turing.FinTM.foldRead_tape}
\label{Turing.FinTM.foldRead_tape}
\leanok
\uses{Turing.FinTM.foldPos, Turing.FinTM.foldPos_nonneg, Turing.FinTM.foldRead, Turing.FinTM.foldSide, Turing.FinTM.foldTape, Turing.FinTM.foldUnpack_pack}
\difficulty{3}
For every tape content $t \colon \bbz \to \Gamma \cup \{\square\}$ over a base
alphabet $\Gamma$ and every integer $z$, reading the folded tape of $t$ at the
physical position $\varphi(z)$ on the side of $z$ returns exactly the original
content $t(z)$, where $\varphi(z) = z$ for $z \ge 0$ and $\varphi(z) = -z-1$
otherwise.
\end{lemma}

\begin{definition}[Writing to the active payload]
\lean{Turing.FinTM.foldWrite}
\label{Turing.FinTM.foldWrite}
\leanok
\uses{Turing.FinTM.FoldSymbol, Turing.FinTM.foldPack, Turing.FinTM.foldUnpack}
Writing a possibly blank base symbol $a$ into a physical cell over the folded
alphabet, relative to a side: unpack the cell, replace the first payload by $a$ when
the side is the nonnegative one and the second payload otherwise, keep the origin
flag and the untouched payload, and repack canonically (so a fully blank untagged
result becomes a physical blank).
\end{definition}

\begin{lemma}[Folded writes commute with virtual updates]
\lean{Turing.FinTM.foldTape_update}
\label{Turing.FinTM.foldTape_update}
\leanok
\uses{Turing.FinTM.foldPos, Turing.FinTM.foldPos_nonneg, Turing.FinTM.foldSide, Turing.FinTM.foldTape, Turing.FinTM.foldUnpack_pack, Turing.FinTM.foldWrite}
\difficulty{5}
Let $\Gamma$ be an alphabet with decidable equality,
$t \colon \bbz \to \Gamma \cup \{\square\}$ a tape content, $z$ an integer, and $a$ a
possibly blank symbol.  Updating the folded tape of $t$ at the physical cell
$\varphi(z)$ by the side-relative write of $a$ (on the side of $z$) yields exactly
the folded tape of the pointwise update of $t$ at $z$ by $a$; here $\varphi(z) = z$
for $z \ge 0$ and $\varphi(z) = -z-1$ otherwise.
\end{lemma}

\begin{definition}[Translation of a machine action]
\lean{Turing.FinTM.foldAction}
\label{Turing.FinTM.foldAction}
\leanok
\uses{Turing.Action, Turing.FinTM.FoldSymbol, Turing.FinTM.foldEmbedding, Turing.FinTM.foldMove, Turing.FinTM.foldUnpack, Turing.FinTM.foldWrite}
Given, for each of $k$ work tapes, the current side of the fold and the physical
cell under the head, this translates an action of a $k$-work-tape machine over a
base alphabet $\Gamma$ with state set $S$ into an action over the folded alphabet
whose states pair a state of $S$ with a side for each tape (or are absent, meaning
halt): writes become side-relative folded writes, moves are translated across the
fold using the origin flags stored in the scanned cells, the emitted output symbol
is embedded, and the successor state is paired with the updated sides.
\end{definition}

\begin{definition}[Folding of a configuration]
\lean{Turing.FinTM.foldCfg}
\label{Turing.FinTM.foldCfg}
\leanok
\uses{Turing.Cfg, Turing.FinTM.FoldSymbol, Turing.FinTM.foldEmbedding, Turing.FinTM.foldPos, Turing.FinTM.foldSide, Turing.FinTM.foldTape}
The folded image of a configuration $c$ of a $k$-work-tape machine over a base
alphabet $\Gamma$ with state set $S$ on an input $x$: a configuration over the
folded alphabet on the symbol-wise embedded input, whose control state is the state
of $c$ paired with the side of each work head, whose work tapes are the folded
tapes of those of $c$, whose work head positions are the folded coordinates
$\varphi$ of those of $c$, and whose input head position and output are transported
along the embedding.
\end{definition}

\begin{lemma}[Folded input symbol]
\lean{Turing.FinTM.foldCfg_input}
\label{Turing.FinTM.foldCfg_input}
\leanok
\uses{Turing.Cfg, Turing.Cfg.inputSymbol, Turing.FinTM.foldCfg, Turing.FinTM.foldEmbedding}
\difficulty{3}
For every configuration $c$ of a $k$-work-tape machine over a base alphabet
$\Gamma$ on an input $x$, the possibly absent input symbol scanned in the folded
configuration of $c$ equals the image under the alphabet embedding of the input
symbol scanned in $c$ (with absence mapped to absence).
\end{lemma}

\begin{lemma}[Origin flags under a folded head]
\lean{Turing.FinTM.foldCfg_origin}
\label{Turing.FinTM.foldCfg_origin}
\leanok
\uses{Turing.Cfg, Turing.Cfg.workTapeSymbols, Turing.FinTM.foldCfg, Turing.FinTM.foldPos, Turing.FinTM.foldPos_nonneg, Turing.FinTM.foldTape, Turing.FinTM.foldUnpack, Turing.FinTM.foldUnpack_pack}
\difficulty{1}
In the folded image of a configuration $c$ of a $k$-work-tape machine over a base
alphabet, for each work tape the origin flag obtained by unpacking the physical
cell under that tape's head equals the truth value of the statement that the folded
head coordinate of that tape is $0$.
\end{lemma}

\begin{lemma}[Folding commutes with applying an action]
\lean{Turing.FinTM.foldCfg_apply}
\label{Turing.FinTM.foldCfg_apply}
\leanok
\uses{Turing.Action, Turing.Action.apply, Turing.Cfg, Turing.Cfg.workTapeSymbols, Turing.FinTM.foldAction, Turing.FinTM.foldCfg, Turing.FinTM.foldCfg_origin, Turing.FinTM.foldDecode, Turing.FinTM.foldMove, Turing.FinTM.foldMove_correct, Turing.FinTM.foldPos, Turing.FinTM.foldSide, Turing.FinTM.foldTape_update, Turing.FinTM.foldUnpack, Turing.moveInputPos}
\difficulty{5}
Let $\Gamma$ be an alphabet with decidable equality, $a$ an action of a
$k$-work-tape machine over $\Gamma$ with state set $S$, and $c$ a configuration of
such a machine on an input $x$.  Applying to the folded configuration of $c$ the
translated action --- formed from each work head's current side and the physical
cells under the folded heads --- yields the folded configuration of the result of
applying $a$ to $c$.
\end{lemma}

\begin{definition}[Decoding of an embedded input symbol]
\lean{Turing.FinTM.foldDecode}
\label{Turing.FinTM.foldDecode}
\leanok
\uses{Turing.FinTM.FoldSymbol}
The partial inverse of the alphabet embedding: a folded triple whose origin flag is
unset, whose first payload is a base symbol $a$, and whose second payload is blank
decodes to $a$; every other folded triple decodes to nothing.
\end{definition}

\begin{definition}[The stationary halting action]
\lean{Turing.FinTM.foldHalt}
\label{Turing.FinTM.foldHalt}
\leanok
\uses{Turing.Action}
The action of a $k$-work-tape machine that halts immediately: it does not move the
input head, writes nothing and issues no move on any work tape, emits no output
symbol, and has no successor state.
\end{definition}

\begin{definition}[The initialization action]
\lean{Turing.FinTM.foldInitAction}
\label{Turing.FinTM.foldInitAction}
\leanok
\uses{Turing.Action, Turing.FinTM.FoldSymbol}
The single initialization action of the folding simulator, parametrized by a target
base state $q$: without moving any head and without emitting output, it
simultaneously writes the origin marker --- the folded triple with flag set and
both payloads blank --- at the current cell of every work tape, and passes control
to the state $q$ with every tape recorded as being on the nonnegative side.
\end{definition}

\begin{definition}[Guarded start action]
\lean{Turing.FinTM.foldStartAction}
\label{Turing.FinTM.foldStartAction}
\leanok
\uses{Turing.Action, Turing.FinTM.FoldSymbol, Turing.FinTM.foldDecode, Turing.FinTM.foldInitAction}
The action taken from the simulator's fresh initial state when scanning a possibly
absent folded input symbol: if the scanned symbol is absent or decodes to a base
symbol, perform the initialization action toward a given base state $q$; if it lies
outside the range of the alphabet embedding, perform the same origin-marking writes
but halt instead of entering $q$.  Thus no simulated transition precedes the
validity check on the first input symbol.
\end{definition}

\begin{lemma}[Start action on embedded inputs]
\lean{Turing.FinTM.foldStartAction_embed}
\label{Turing.FinTM.foldStartAction_embed}
\leanok
\uses{Turing.FinTM.foldEmbedding, Turing.FinTM.foldInitAction, Turing.FinTM.foldStartAction}
\difficulty{1}
For every base state $q$ and every possibly absent base symbol, the guarded start
action evaluated at the image of that symbol under the alphabet embedding equals
the plain initialization action toward $q$: embedded symbols never trigger the
halting guard.
\end{lemma}

\begin{definition}[The folding simulator machine]
\lean{Turing.FinTM.foldTM}
\label{Turing.FinTM.foldTM}
\leanok
\uses{Turing.FinTM, Turing.FinTM.FoldSymbol, Turing.FinTM.foldAction, Turing.FinTM.foldDecode, Turing.FinTM.foldHalt, Turing.FinTM.foldRead, Turing.FinTM.foldStartAction}
For a multi-tape machine $M$ over a finite alphabet $\Gamma$ with decidable
equality, the simulating machine over the folded alphabet: it has the same number
of work tapes as $M$, and its states are either a fresh initial state or a state of
$M$ paired with a side for each work tape.  From the fresh state it performs the
guarded initialization; thereafter it decodes the scanned input symbol (executing
the stationary halt if the symbol lies outside the embedding's range), reads the
active payload of each work tape, and executes the translation of the corresponding
transition of $M$.
\end{definition}

\begin{lemma}[One-step simulation commutation]
\lean{Turing.FinTM.foldCfg_step}
\label{Turing.FinTM.foldCfg_step}
\leanok
\uses{Turing.Cfg, Turing.Cfg.inputSymbol, Turing.Cfg.workTapeSymbols, Turing.FinTM, Turing.FinTM.foldCfg, Turing.FinTM.foldCfg_apply, Turing.FinTM.foldCfg_input, Turing.FinTM.foldDecode, Turing.FinTM.foldEmbedding, Turing.FinTM.foldRead, Turing.FinTM.foldRead_tape, Turing.FinTM.foldSide, Turing.FinTM.foldTM, Turing.MultiTapeTM.step}
\difficulty{4}
Let $M$ be a multi-tape machine over a finite alphabet $\Gamma$ with decidable
equality and let $c$ be a configuration of $M$ on an input $x$ over $\Gamma$.  One
step of the folding simulator of $M$ on the folded configuration of $c$ equals the
folded configuration of one step of $M$ on $c$.
\end{lemma}

\begin{lemma}[Initialization reaches the folded start configuration]
\lean{Turing.FinTM.foldCfg_init}
\label{Turing.FinTM.foldCfg_init}
\leanok
\uses{Turing.Action.apply, Turing.Cfg.inputSymbol, Turing.FinTM, Turing.FinTM.foldCfg, Turing.FinTM.foldCfg_input, Turing.FinTM.foldEmbedding, Turing.FinTM.foldInitAction, Turing.FinTM.foldPack, Turing.FinTM.foldPos, Turing.FinTM.foldStartAction, Turing.FinTM.foldStartAction_embed, Turing.FinTM.foldTM, Turing.FinTM.foldTape, Turing.MultiTapeTM.initCfg, Turing.MultiTapeTM.step}
\difficulty{4}
Let $M$ be a multi-tape machine over a finite alphabet $\Gamma$ with decidable
equality and let $x$ be an input over $\Gamma$.  One step of the folding simulator
from its initial configuration on the symbol-wise embedded input $x$ yields exactly
the folded configuration of the initial configuration of $M$ on $x$.
\end{lemma}

\begin{lemma}[Lockstep simulation with a one-step offset]
\lean{Turing.FinTM.foldCfg_run}
\label{Turing.FinTM.foldCfg_run}
\leanok
\uses{Turing.FinTM, Turing.FinTM.foldCfg, Turing.FinTM.foldCfg_init, Turing.FinTM.foldCfg_step, Turing.FinTM.foldEmbedding, Turing.FinTM.foldTM, Turing.MultiTapeTM.initCfg, Turing.MultiTapeTM.runFrom, Turing.MultiTapeTM.runFrom_succ_eq_step, Turing.MultiTapeTM.runFrom_succ_eq_step', Turing.MultiTapeTM.runFrom_zero}
\difficulty{4}
Let $M$ be a multi-tape machine over a finite alphabet $\Gamma$ with decidable
equality, $x$ an input over $\Gamma$, and $t \in \bbn$.  Running the folding
simulator for $t+1$ steps from its initial configuration on the symbol-wise
embedded input equals the folded configuration of running $M$ for $t$ steps from
its initial configuration on $x$.
\end{lemma}

\begin{lemma}[Writes preserve the origin flag]
\lean{Turing.FinTM.foldWrite_origin}
\label{Turing.FinTM.foldWrite_origin}
\leanok
\uses{Turing.FinTM.FoldSymbol, Turing.FinTM.foldUnpack, Turing.FinTM.foldUnpack_pack, Turing.FinTM.foldWrite}
\difficulty{1}
For every side $s$, every physical cell content $w$ over the folded alphabet, and
every possibly blank base symbol $a$, the origin flag obtained by unpacking the
result of the side-relative write of $a$ into $w$ equals the origin flag obtained
by unpacking $w$.
\end{lemma}

\begin{lemma}[Translated moves stay nonnegative]
\lean{Turing.FinTM.foldMove_nonneg}
\label{Turing.FinTM.foldMove_nonneg}
\leanok
\uses{Turing.FinTM.foldMove}
\difficulty{3}
Let $p$ be a nonnegative integer, $s$ a side, and $d \in \{-1,0,+1\}$ a move
direction.  Then $p$ plus the physical displacement of the translated move computed
from $s$, the flag ``$p = 0$'', and $d$ is still nonnegative.
\end{lemma}

\begin{definition}[The safety invariant]
\lean{Turing.FinTM.foldSafe}
\label{Turing.FinTM.foldSafe}
\leanok
\uses{Turing.Cfg, Turing.FinTM.FoldSymbol, Turing.FinTM.foldUnpack}
The safety invariant on a configuration $c$ of a $k$-work-tape machine over the
folded alphabet whose states are either a fresh initial state or a base state
paired with per-tape sides, on an arbitrary folded-alphabet input: the control
state of $c$ is not the fresh initial state (it may be halted), every work head
position is nonnegative, and on every work tape the origin flag obtained by
unpacking each nonnegative cell $p$ is set exactly when $p = 0$.
\end{definition}

\begin{lemma}[Translated actions preserve safety]
\lean{Turing.FinTM.foldAction_safe}
\label{Turing.FinTM.foldAction_safe}
\leanok
\uses{Turing.Action, Turing.Action.apply, Turing.Cfg, Turing.Cfg.workTapeSymbols, Turing.FinTM.FoldSymbol, Turing.FinTM.foldAction, Turing.FinTM.foldMove, Turing.FinTM.foldMove_nonneg, Turing.FinTM.foldSafe, Turing.FinTM.foldUnpack, Turing.FinTM.foldWrite_origin}
\difficulty{4}
Let $c$ be a configuration of a $k$-work-tape machine over the folded alphabet, on
an arbitrary folded-alphabet input, satisfying the safety invariant (control not in
the fresh initial state, nonnegative heads, correct origin flags at all nonnegative
cells).  Then for every assignment of sides to the work tapes and every base-machine
action $a$, applying the translation of $a$ --- relative to those sides and the
cells under the heads of $c$ --- again yields a configuration satisfying the safety
invariant.
\end{lemma}

\begin{lemma}[The halting action preserves safety]
\lean{Turing.FinTM.foldHalt_safe}
\label{Turing.FinTM.foldHalt_safe}
\leanok
\uses{Turing.Action.apply, Turing.Cfg, Turing.FinTM.FoldSymbol, Turing.FinTM.foldHalt, Turing.FinTM.foldSafe}
\difficulty{1}
If a configuration of a $k$-work-tape machine over the folded alphabet, on an
arbitrary folded-alphabet input, satisfies the safety invariant, then so does the
configuration obtained from it by applying the stationary halting action.
\end{lemma}

\begin{lemma}[Simulator steps preserve safety]
\lean{Turing.FinTM.foldSafe_step}
\label{Turing.FinTM.foldSafe_step}
\leanok
\uses{Turing.Cfg, Turing.Cfg.inputSymbol, Turing.FinTM, Turing.FinTM.FoldSymbol, Turing.FinTM.foldAction_safe, Turing.FinTM.foldDecode, Turing.FinTM.foldHalt_safe, Turing.FinTM.foldSafe, Turing.FinTM.foldTM, Turing.MultiTapeTM.step}
\difficulty{4}
Let $M$ be a multi-tape machine over a finite alphabet with decidable equality and
let $c$ be a configuration of the folding simulator of $M$, on an arbitrary
folded-alphabet input, satisfying the safety invariant.  Then the configuration
obtained after one step of the simulator also satisfies the safety invariant.
\end{lemma}

\begin{lemma}[Safety after the first step]
\lean{Turing.FinTM.foldSafe_init}
\label{Turing.FinTM.foldSafe_init}
\leanok
\uses{Turing.Action.apply, Turing.FinTM, Turing.FinTM.FoldSymbol, Turing.FinTM.foldInitAction, Turing.FinTM.foldSafe, Turing.FinTM.foldStartAction, Turing.FinTM.foldTM, Turing.FinTM.foldUnpack, Turing.MultiTapeTM.initCfg, Turing.MultiTapeTM.step}
\difficulty{4}
Let $M$ be a multi-tape machine over a finite alphabet with decidable equality.
For every input over the folded alphabet --- including inputs containing symbols
outside the range of the alphabet embedding --- the configuration reached after one
step of the folding simulator from its initial configuration satisfies the safety
invariant.
\end{lemma}

\begin{lemma}[The simulator uses its tapes unidirectionally]
\lean{Turing.FinTM.foldTM_nonnegative}
\label{Turing.FinTM.foldTM_nonnegative}
\leanok
\uses{Turing.FinTM, Turing.FinTM.NonnegativeHeads, Turing.FinTM.foldSafe, Turing.FinTM.foldSafe_init, Turing.FinTM.foldSafe_step, Turing.FinTM.foldTM, Turing.MultiTapeTM.initCfg, Turing.MultiTapeTM.runFrom, Turing.MultiTapeTM.runFrom_succ_eq_step, Turing.MultiTapeTM.runFrom_succ_eq_step', Turing.MultiTapeTM.step}
\difficulty{4}
For every multi-tape machine $M$ over a finite alphabet with decidable equality,
the folding simulator of $M$ has nonnegative heads: for every input over the folded
alphabet, every number of steps, and every work tape, the position of that tape's
head in the corresponding initialized run is nonnegative.
\end{lemma}

\begin{theorem}[Unidirectional tapes suffice]
\lean{Turing.FinTM.nonnegative_heads}
\label{Turing.FinTM.nonnegative_heads}
\leanok
\uses{Turing.FinTM, Turing.FinTM.ComputesFunInTime, Turing.FinTM.ComputesFunInTimeVia, Turing.FinTM.FoldSymbol, Turing.FinTM.NonnegativeHeads, Turing.FinTM.computesInTime_iff, Turing.FinTM.foldCfg, Turing.FinTM.foldCfg_run, Turing.FinTM.foldEmbedding, Turing.FinTM.foldTM, Turing.FinTM.foldTM_nonnegative}
\difficulty{3}
Let $\Gamma$ be a finite alphabet with decidable equality, let $M$ be a multi-tape
Turing machine over $\Gamma$, and suppose $M$ computes a function
$f \colon \Gamma^* \to \Gamma^*$ within a time bound $T \colon \bbn \to \bbn$.
Then there exist a finite alphabet $\Gamma'$ with decidable equality, an injection
$e \colon \Gamma \hookrightarrow \Gamma'$, a machine $M'$ over $\Gamma'$, and a
constant $c \in \bbn$ such that: in every initialized run of $M'$ no work head ever
visits a negative cell, $M'$ has the same number of work tapes as $M$, and $M'$
computes $f$ through the embedding $e$ within the time bound
$n \mapsto c \cdot (T(n) + 1)$.
\end{theorem}
